\section{ANTLR dan Perbandingan Pendekatan}
\label{sec:ch13-antlr}

\subsection{ANTLR}

ANTLR adalah parser generator yang menghasilkan parser LL(*) dengan
dukungan \eng{visitor} dan \eng{listener} \citep{antlr4}. Grammar ditulis dalam
berkas \file{.g4}. Contoh ringkas:
\begin{lstlisting}[language=Java]
grammar Expr;
prog   : stat+ ;
stat   : expr NEWLINE ;
expr   : expr ('*'|'/') expr
       | expr ('+'|'-') expr
       | INT
       | '(' expr ')'
       ;
INT    : [0-9]+ ;
NEWLINE: '\r'? '\n' ;
WS     : [ \t]+ -> skip ;
\end{lstlisting}

ANTLR membangun pohon parse otomatis. Pengguna dapat menulis
\eng{visitor} untuk mengubah pohon menjadi AST atau melakukan
analisis. ANTLR juga menyediakan pesan kesalahan dan pemulihan yang
cukup baik.

\subsection{Perbandingan}

\begin{table}[H]
\centering
\caption{Perbandingan Bison, ANTLR, dan parser tangan}
\label{tab:ch13-perbandingan}
\begin{tabularx}{\textwidth}{@{}lXXX@{}}
\toprule
\textbf{Aspek} & \textbf{Bison} & \textbf{ANTLR} & \textbf{Tangan} \\
\midrule
Algoritma & LALR(1) & LL(*) & Recursive descent \\
Bahasa aksi & C/C++ & Java, C++, Python & Bahasa pilihan \\
Pohon & Dibangun manual & Otomatis & Dibangun manual \\
Kontrol & Tinggi & Sedang & Penuh \\
Diagnostik & Perlu kode tambahan & Bawaan baik & Penuh, tetapi
manual \\
Kurva belajar & Sedang & Rendah--sedang & Rendah \\
\bottomrule
\end{tabularx}
\end{table}

\subsection{Kapan Memilih}

Bison cocok untuk proyek C/C++ yang membutuhkan kontrol dan kinerja.
ANTLR cocok untuk prototipe cepat, bahasa dengan grammar kompleks,
dan tim yang bekerja di ekosistem Java/Python. Parser tangan cocok
untuk bahasa kecil, pembelajaran, dan kebutuhan diagnostik yang
sangat spesifik.

\begin{infobox}
Pilihan alat bukan penentu keberhasilan. Yang lebih penting adalah
grammar yang benar, pengujian yang memadai, dan pemahaman tentang
konflik serta pemulihan kesalahan.
\end{infobox}
